\documentclass[11pt]{article}

\usepackage[T1]{fontenc}
\usepackage[utf8]{inputenc}
\usepackage[margin=1in]{geometry}
\usepackage{amsmath,amssymb}
\usepackage{graphicx}
\usepackage{booktabs}
\usepackage{hyperref}

\title{Paper Title}
\author{Author Name\\
  \small Institution\\
  \small \texttt{author@example.com}}
\date{\today}

\begin{document}

\maketitle

\begin{abstract}
Write a concise summary of the paper: its purpose, approach, main results, and
conclusions.
\end{abstract}

\noindent\textbf{Keywords:} keyword one, keyword two, keyword three

\section{Introduction}
Introduce the topic, explain its significance, and state the research question
or objective. Summarize relevant prior work and cite it with
\verb|\cite{example}|.

\section{Methods}
Describe the materials, data, and methods used so that the work can be
understood and reproduced. For example, an equation can be written as
\begin{equation}
  y = mx + b.
\end{equation}

\section{Results}
Present the findings clearly. Tables can be formatted with \texttt{booktabs}:
\begin{table}[ht]
  \centering
  \caption{Example results table. Replace the sample values with your data.}
  \label{tab:example}
  \begin{tabular}{lc}
    \toprule
    Measure & Value \\
    \midrule
    Example measure & 0.0 \\
    \bottomrule
  \end{tabular}
\end{table}

\section{Discussion}
Interpret the results, relate them to prior work, and discuss limitations.

\section{Conclusion}
Summarize the main findings and their implications. Note possible directions
for future work.

\begin{thebibliography}{9}
\bibitem{example}
A. Author, ``Title of the cited work,'' Publication, year.
\end{thebibliography}

\end{document}
